\section{Spectral preliminaries and relative singularity}
\label{sec:spectral}

This section and Section~\ref{sec:rigidity} establish the probability-law
rigidity used in the coloring transfer of Section~\ref{sec:transfer}.
The laws live on the character group of the discrete algebraic plane.
We state the rigidity theorem after defining that group, then prove
that exclusion of ordinary continuous characters persists in the
conditional difference laws arising along compact towers.  This
relative singularity is the input to the Fourier argument in
Section~\ref{sec:rigidity}.

\subsection{Algebraic rotations and continuous characters}

\begin{definition}\label{def:spectral-groups}
Identify the plane with $\C$, and set
\[
 F=\overline{\Q}\cap\R,\qquad E=F(i),\qquad
 K=\{u\in E:u\bar u=1\}.
\]
We regard $E$ as an additive discrete group and $K$ as a multiplicative
discrete group. Let
\[
 D=\widehat{E_{\mathrm{disc}}}
\]
be the compact character group of $E$. A point $d\in D$ is a homomorphism
from the additive group $E$ to the unit complex group $\T$. We use additive
notation on $D$, so $(d+d')(z)=d(z)d'(z)$. For $a,z\in E$, put
\[
 (ad)(z)=d(az),\qquad \psi_z(d)=d(z).
\]
The subgroup of ordinary continuous characters is
\begin{equation}\label{eq:spectral-continuous-characters}
 C=j(\R^2),\qquad
 j(\xi)(z)=\exp(i\xi\cdot z),
\end{equation}
where the dot product uses the real coordinates of $z$.
\end{definition}

Both $E$ and $K$ are countable abelian groups. Every nonzero power map
$u\mapsto u^n$ on $K$ is onto: a complex $n$th root of an algebraic number
of modulus one is again algebraic and has modulus one; negative powers
also cause no difficulty. Thus $K$ is divisible. The compact group $D$ is
metrizable because $E$ is countable.

The map $j$ is continuous and injective. Indeed, a continuous plane
character that is one on the dense subgroup $E$ is one everywhere, and
its frequency is zero. Its image is the union of the compact images of
closed balls in $\R^2$, so $C$ is Borel. The inverse of $j$ on each such
compact image is continuous; consequently $j^{-1}:C\to\R^2$ is Borel.
For $a\ne0$, multiplication by $a$ is an automorphism of $D$ and carries
$C$ onto itself. On the frequency parameters it acts by the transpose
of the real linear map of multiplication by $a$. In particular, an
element of $K$ acts on these parameters by an orthogonal map.

\begin{definition}\label{def:spectral-wild}
A $K$-invariant Borel probability measure $\nu$ on $D$ is called
\emph{wild} if $\nu(C)=0$. Its Fourier coefficients are denoted by
\[
 f_\nu(z)=\int_D\psi_z(d)\,d\nu(d),\qquad z\in E.
\]
\end{definition}

\begin{theorem}[Rigidity of wild character laws]\label{thm:rigidity}
Every $K$-invariant probability measure $\nu$ on $D$ with $\nu(C)=0$
is Haar measure. Equivalently, writing $m_D$ for normalized Haar measure,
\begin{equation}\label{eq:source-3}
 \nu\text{ is wild}\quad\Longrightarrow\quad\nu=m_D.
\end{equation}
\end{theorem}

The Fourier coefficients of Haar probability vanish at every nonzero
$z\in E$.  This is the conclusion needed in Section~\ref{sec:transfer}:
a rotation-fixed label indicator has a $K$-invariant finite spectral
measure, and its part off $C$ is again invariant.  If that part is
nonzero, normalization allows the theorem to be applied; its correlations
at every nonzero algebraic translation then vanish.  The factor argument in
Lemma~\ref{lem:palettes-continuous-factor} will make the $C$-supported
projections nonnegative, with $L^2$-continuous translation orbits.

There is a classical analogue for measurable radial positive-definite
functions on the full real plane: each is the sum of a continuous
positive-definite function and a nonnegative multiple of
$\ind_{\{0\}}$ \citep{Crum1956,GneitingSasvari1999}.
Here the Fourier coefficients are initially defined only on $E$, with
invariance under its algebraic rotations $K$; no measurable radial
positive-definite extension to $\R^2$ is assumed.
We prove the required rigidity directly on $D$.

The proof of the theorem is completed in Section~\ref{sec:rigidity}.
Here we first obtain null-event and line-average consequences of
$\nu(C)=0$.  We then show that this condition survives when the law is
replaced by the difference of two conditionally independent samples
over a factor in a compact tower.

\begin{lemma}\label{lem:spectral-null-events}
Let $\nu$ be wild. Then every coset of $C$ has $\nu$-measure zero, and
the law of the difference of two independent $\nu$-samples gives $C$
measure zero. Moreover, for each $a\in E\setminus\{0\}$, the event
\[
 L_a=\{d\in D:t\mapsto d(at)\text{ is continuous on }F
                 \text{ in its usual topology}\}
\]
is Borel and has $\nu$-measure zero.
\end{lemma}

\begin{proof}
If a coset $d+C$ has positive mass, its $K$-orbit is finite, since the
cosets in that orbit are disjoint and have the same mass. The action on
this finite orbit gives a finite image of the divisible abelian group
$K$. Such an image is trivial: a finite abelian group cannot be divisible
unless it is trivial. Hence every $u\in K$ fixes the coset, and
$(u-1)d\in C$. Taking $u\ne1$ and applying the inverse scalar shows
$d\in C$, contrary to $\nu(C)=0$. The difference assertion now follows
by integrating $\nu(d+C)=0$ against $\nu(d)$; the relation
$d'-d\in C$ is Borel.

Restriction to the line $aF$ is a continuous map from $D$ to
$\widehat{F_{\mathrm{disc}}}$. The ordinary continuous characters on $F$
are exactly $t\mapsto e^{i\lambda t}$, $\lambda\in\R$. Their
parametrization is continuous and injective, with Borel image by compact
exhaustion, as for $j$. This proves the Borel assertion for $L_a$.

If $aF,bF$ are nonparallel, then $a,b$ form an $F$-basis of $E$. A
character continuous on both lines has the form
\[
 d(at+bs)=e^{i(\lambda t+\eta s)}\qquad(t,s\in F)
\]
for real $\lambda,\eta$. The corresponding real coordinate change is
invertible and continuous, so this character belongs to $C$. Thus
$L_a\cap L_b\subset C$ for nonparallel lines. Rotations give infinitely
many distinct unoriented lines $uaF$, and their events have the same
$\nu$-measure. Any finite collection of them is pairwise disjoint modulo
$\nu$-null sets. Their common measure must therefore be zero.
\end{proof}

\subsection{A line average without a measurability assumption}

Fix a translation-invariant mean $\mathfrak m$ on the bounded functions
on the discrete abelian group $F/\Z$. Such a normalized positive mean
exists by amenability of abelian groups; see
\citet[Theorem~G.2.1]{Bekka2008}. Because the group is discrete,
the mean is defined on every bounded function. Invariant means and
the finite-set averaging criterion originate in the work of
\citet{vonNeumann1929} and \citet{Folner1955}. The restriction of
$\mathfrak m$ to ordinary continuous periodic functions is integration
over the usual circle. To check this, every nonconstant ordinary circle
character is multiplied by a nontrivial scalar under some translation
from $F/\Z$, and hence has mean zero. Trigonometric-polynomial
approximation gives the assertion for all continuous periodic functions.

For a bounded function $h:F\to\C$ supported in a bounded interval, define
\begin{equation}\label{eq:spectral-invariant-integral}
 I(h)=\mathfrak m\left(t+\Z\longmapsto
                         \sum_{k\in\Z}h(t+k)\right).
\end{equation}
The periodization is well defined and bounded: only a uniformly bounded
number of summands can be nonzero. The functional $I$ is positive and
translation invariant under $F$. For a continuous compactly supported
function $\phi$ on $\R$, it satisfies
\[
 I(\phi|_F)=\int_\R\phi(t)\,dt.
\]
We also use $I$ to integrate bounded, boundedly supported functions with
values in a Hilbert space, in the weak sense. Scalar pairings define a
bounded functional on the Hilbert space, and Hilbert-space duality gives
the resulting vector. In particular, if $V(t)$ is such a function, then
\[
 \norm{I(V)}\le I\bigl(\norm{V(t)}\bigr).
\]
This construction does not impose ordinary real-variable measurability
on $V$ or on a Fourier coefficient restricted to $F$.

\begin{lemma}[Vanishing line average]\label{lem:line-average}
If $\nu$ is wild, then
\begin{equation}\label{eq:source-2}
 I\bigl(\phi(t)f_\nu(at)\bigr)=0
 \qquad\bigl(a\in E\setminus\{0\},\ \phi\in C_c(\R)\bigr).
\end{equation}
\end{lemma}

\begin{proof}
On $L^2(\nu)$ consider the unitary representation of the additive group
$F$ given by
\[
 T_tv=\psi_{at}v.
\]
First, this representation has no nonzero vector that is strongly
continuous for the usual topology on $F$. Suppose that $v$ is such a
vector. Put
\[
 c_v(t)=\int_D|v|^2\psi_{at}\,d\nu.
\]
It is positive definite and uniformly continuous on $F$, since
\[
 |c_v(t+s)-c_v(t)|\le\norm{v}_2\norm{T_sv-v}_2.
\]
It therefore extends continuously to $\R$. The extension remains
positive definite, by approximation of each finite set of real
arguments by elements of $F$.

Apply the ordinary Bochner theorem to this continuous positive-definite
function on $\R$; see \citet[Theorem~D.2.2]{Bekka2008}.
It is important here that the coefficient has first been extended to
$\R$; no local compactness of $F$ in its usual topology is assumed.
Push its finite positive
representing measure to $\widehat{F_{\mathrm{disc}}}$ using the
continuous-character parametrization. Its Fourier coefficients agree
with those of the restriction-map pushforward of $|v|^2\nu$. Uniqueness
of Fourier coefficients for finite measures on the compact character
group identifies these two measures. Indeed, the characters form a
unital conjugation-closed algebra separating its points, so their
uniform density in the continuous functions gives this uniqueness. The pushforward of $|v|^2\nu$ is
thus concentrated on the ordinary continuous characters, so $|v|^2\nu$
is concentrated on $L_a$. Lemma~\ref{lem:spectral-null-events} gives
$\nu(L_a)=0$, and hence $v=0$.

Now form the weak integral
\[
 V_\phi=I\bigl(\phi(t)T_t1\bigr)\in L^2(\nu).
\]
Translation invariance of $I$ gives
\[
 T_sV_\phi-V_\phi
 =I\bigl((\phi(t-s)-\phi(t))T_t1\bigr),
\]
and consequently
\[
 \norm{T_sV_\phi-V_\phi}_2
 \le I(|\phi(\cdot-s)-\phi|)
 =\int_\R|\phi(t-s)-\phi(t)|\,dt\longrightarrow0
\]
as $s\in F$ tends to zero. Thus $V_\phi$ is strongly continuous and
must be zero by the preceding argument. Pairing with the constant
function proves \eqref{eq:source-2}.
\end{proof}

\subsection{A measurable equivariant center}

The subgroup $C$ need not be closed. We will nevertheless choose a
canonical representative of a positive-mass $C$-coset when a probability
kernel on that coset is available. The following elementary location
rule avoids any moment assumption on that kernel.

\begin{lemma}\label{lem:spectral-location}
For every Borel probability measure $\lambda$ on $\R^2$, the function
\begin{equation}\label{eq:spectral-location-objective}
 J_\lambda(s)=\int_{\R^2}
 \left(\sqrt{1+|s-t|^2}-\sqrt{1+|t|^2}\right)\,d\lambda(t)
\end{equation}
has a unique minimizer $L(\lambda)$. The map $\lambda\mapsto L(\lambda)$
is Borel for the usual Borel structure on the space of probability
measures. It is equivariant under translations and orthogonal maps.
\end{lemma}

\begin{proof}
The difference in \eqref{eq:spectral-location-objective} is bounded in
absolute value by $|s|$. The objective is therefore finite without a
first-moment assumption. It is Lipschitz in $s$ and strictly convex,
because $s\mapsto\sqrt{1+|s-t|^2}$ is strictly convex for every $t$.

Choose $R>0$ with $p=\lambda(\overline{B(0,R)})>1/2$. On this ball the
integrand is at least $|s|-R-\sqrt{1+R^2}$, and elsewhere it is at least
$-|s|$. Hence
\[
 J_\lambda(s)\ge
 (2p-1)|s|-p\bigl(R+\sqrt{1+R^2}\bigr)\longrightarrow\infty.
\]
Coercivity and continuity give a minimizer, and strict convexity gives
uniqueness.

For each fixed $s$, the integrand is bounded and continuous in $t$, so
$\lambda\mapsto J_\lambda(s)$ is Borel. The infimum over rational
points equals the minimum. Enumerate $\Q^2$ and, for each positive
integer $n$, choose its first point whose objective is less than that
infimum plus $1/n$. These choices are measurable. For a fixed $\lambda$
they converge to the unique minimizer: coercivity confines all sufficiently
accurate choices to a compact set, and every cluster point minimizes
$J_\lambda$. Their limit is therefore the Borel map $L$.

If $\lambda'$ is the law of $t-r$ for $t$ of law $\lambda$, then an
identity of bounded differences gives
\[
 J_{\lambda'}(s)=J_\lambda(s+r)-J_\lambda(r).
\]
Thus $L(\lambda')=L(\lambda)-r$. Orthogonal covariance follows from
$J_{O_*\lambda}(Os)=J_\lambda(s)$. These identities prove both asserted
equivariances.
\end{proof}

\subsection{Compact towers and relative singularity}

We use standard probability spaces and probability-preserving actions
of the countable group $K$. Factors are countably generated modulo null
sets. For an action write $U_kH(x)=H(kx)$, and identify functions on a
factor with their pullbacks. Conditional kernels, factor actions, and
equivariant maps may be chosen on common conull sets whenever needed;
the acting group is countable.

We use disintegration in the following precise form. A factor map
$\pi:X\to Y$ between standard Borel probability models admits a
measurable probability kernel $y\mapsto\mu_y$, concentrated on
$\pi^{-1}(y)$ for almost every $y$, such that
$\mu=\int\mu_y\,d\nu(y)$. The kernel gives conditional expectations
by integration and is unique almost everywhere; see
\citet[Theorem~4 and Remark~5]{TaoDisintegration2008}.
For an equivariant factor map, uniqueness makes the kernel equivariant
for each group element. Intersecting the corresponding conull sets
makes equivariance simultaneous for the countable group $K$.
This is the fiber-measure input to all relative products below.

Consider a tower $(Y_\alpha)_{\alpha\le\theta}$ of factors of an
ergodic system $X$, where $\theta$ is a countable ordinal. Start with the
trivial factor $Y_0$, and at a limit ordinal take the factor generated
by the preceding ones. Require the following approximation property at
each successor $Y$ over its predecessor $W$: for every bounded $H$ on
$Y$ and every $\varepsilon>0$, there are bounded functions
$H_1,\ldots,H_q$ on $Y$ and a finite constant $B$ such that, for every
$k\in K$,
\begin{equation}\label{eq:source-4}
 \left\|U_kH-\sum_{r=1}^q b_r^k H_r\right\|_2<\varepsilon
 \quad\text{for some }b_r^k\in L^\infty(W)
 \text{ with }\|b_r^k\|_\infty\le B.
\end{equation}
The basis functions and bounds are fixed before $k$ is chosen.

\begin{lemma}[Relative singularity]\label{lem:relative-singularity}
Let $X$ be an ergodic standard probability-preserving $K$-system, and let
$d:X\to D$ be an equivariant measurable map with wild law. For every
countable tower of factors as above, and every factor $Y$ in that tower,
the law $\nu_Y$ of
\[
 d(x)-d(x')
\]
for two conditionally independent samples of $X$ over $Y$ satisfies
$\nu_Y(C)=0$.
\end{lemma}

\begin{proof}
All relative difference laws in the statement are $K$-invariant, by
equivariance of conditional kernels. We use transfinite induction in
its universal form: at each ordinal the assertion ranges over
\emph{all} ambient ergodic standard systems, equivariant maps with wild
law, and towers satisfying \eqref{eq:source-4}. This allows an earlier
induction assertion to be applied below to a new factor system and a
new equivariant map.

At the trivial factor the assertion is
Lemma~\ref{lem:spectral-null-events}. Suppose that it holds at every
earlier ordinal, and that it fails at a noninitial stage $Y$ of a tower
in a system $X$. Write $\kappa_y$ for the conditional law of $d$ over
$y\in Y$.

\paragraph{Positive-mass cosets and their centers.}
The function
\[
 r(y,z)=\kappa_y(z+C)
       =\int_D\ind_C(d'-z)\,d\kappa_y(d')
\]
is Borel in $(y,z)$, by Borel parameter integration. This uses only that
$C$ is Borel, not that it is closed. The function
$x\mapsto r(y(x),d(x))$ is invariant and has integral $\nu_Y(C)>0$.
Ergodicity makes it a positive constant $a$ almost everywhere.

For almost every $y$, therefore, $\kappa_y$-almost every sampled coset
has mass $a$. Such cosets are disjoint and there are only finitely many;
they exhaust the probability. Their number $m$ satisfies $ma=1$, and
is consequently independent of $y$. Thus $\kappa_y$ is concentrated
on $m$ distinct cosets of mass $1/m$ each. None is $C$ almost surely,
because the original marginal law gives $C$ measure zero.

On the Borel set $r(y,z)>0$, define the probability kernel
$\lambda_{y,z}$ on $\R^2$ by
\[
 \lambda_{y,z}(A)=\frac{1}{r(y,z)}
 \int_D\ind_C(d'-z)\ind_A\bigl(j^{-1}(d'-z)\bigr)
 \,d\kappa_y(d')
\]
for Borel $A\subset\R^2$, interpreting the integrand as zero off $C$.
The Borel inverse of $j$ makes this a measurable kernel. Set
\[
 b(y,z)=z+j\bigl(L(\lambda_{y,z})\bigr).
\]
If $z$ is replaced by $z+j(r_0)$, the kernel is translated by $-r_0$;
Lemma~\ref{lem:spectral-location} shows that $b(y,z)$ is unchanged.
It is therefore constant on each positive-mass coset, and remains in
that coset. The same lemma and the orthogonal action of $K$ on frequency
parameters show
\[
 b(ky,kz)=k b(y,z)
\]
on the relevant conull sets.

Let $Z$ be the factor of $X$ generated by $Y$ and
$b(x)=b(y(x),d(x))$. It is a standard ergodic factor. The map
$b:Z\to D$ is equivariant and has wild law. Its conditional law
$\beta_y$ over $Y$ is uniform on the $m$ distinct centers of the
positive-mass cosets. We have converted the hypothetical failure of
singularity into a finite set of equivariant centers over each fiber.
The earlier induction hypothesis and the line average will give
selected real algebraic frequencies with small conditional collision
probabilities on a common fiber.  Compact approximation, however, will
force many pairs of unit-direction phase lists to match one pair of
lists.  Writing each selected frequency as a sum of its two directions
then places its phases near at most $m^2$ products, forcing a larger
collision probability.

\paragraph{A finite list of phases.}
For each $y$, consider the unordered list of $m$ phases $\psi_1$ at
the points of $\beta_y$, retaining repetitions when phases coincide.
Encode this list by the vector $H(y)\in\C^m$ of its elementary
symmetric functions. The vector is measurable: its power sums are
$m\int\psi_1^n\,d\beta_y$, and Newton's identities recover the
elementary symmetric functions. It is bounded.

The encoding is a continuous injection of the compact space
$\T^m/S_m$ into $\C^m$. On unordered lists use maximum matching
distance,
\[
 d_{\mathrm{match}}([z_1,\ldots,z_m],[z'_1,\ldots,z'_m])
 =\min_{\pi\in S_m}\max_{1\le r\le m}|z_r-z'_{\pi(r)}|.
\]
The inverse encoding is uniformly continuous for this metric.
Equivariance of $\beta_y$ and $\psi_1(kd)=\psi_k(d)$ show that
$U_kH$ encodes the list of phases at $\psi_k$.

Choose $\rho>0$ so small that uniform-circle probability of
$\{|z-1|\le5\rho\}$ is less than $1/(16m^2)$. There are an integer
$N\ge1$ and a constant $c>0$ with the following property: if a
probability $\omega$ on $\T$ satisfies
\[
 \left|\int z^n\,d\omega(z)\right|<c\qquad(1\le n\le N),
\]
then
\begin{equation}\label{eq:spectral-small-arc}
 \omega\{|z-1|\le4\rho\}<\frac1{4m^2}.
\end{equation}
Indeed, choose a continuous majorant of the smaller arc that takes
values in $[0,1]$ and vanishes outside the larger arc. Approximate it
uniformly within $1/(64m^2)$ by a real trigonometric polynomial of
degree at most $N$. The constant coefficient is within that error of
its circle integral. Choose $c$ so that $c$ times the sum of the
absolute values of its nonconstant coefficients is less than
$1/(32m^2)$. The integral against $\omega$ of the majorant is then
less than $1/(8m^2)$, which proves
\eqref{eq:spectral-small-arc}. Negative moments are controlled by the
positive ones because $\omega$ is a positive measure.

Choose $s>0$ such that two encoding vectors at distance less than $s$
have phase lists matching within $\rho$. Put $\tau=s/100$, and fix
$\gamma>0$ such that
\begin{equation}\label{eq:spectral-error-choice}
 \frac{2\gamma^2}{\tau^2}<10^{-4}.
\end{equation}

\paragraph{Return to an earlier factor.}
The rotations of $H$ admit vector-valued approximations of the form
\begin{equation}\label{eq:spectral-vector-approximation}
 \left\|U_kH-\sum_{r=1}^q b_r^k H_r\right\|_2<\gamma,
 \qquad \|b_r^k\|_\infty\le B,
\end{equation}
with fixed bounded $\C^m$-valued functions $H_r$ on $Y$ and
coefficients over some earlier factor $W$. At a successor this follows
from \eqref{eq:source-4}, applied to the finitely many coordinates and
combining their data. At a limit, first approximate $H$ in $L^2$,
within $\gamma/2$, by bounded functions on an earlier successor
factor. Such approximation follows from generation of the limit
factor, for example by conditional expectations. Its error remains
unchanged under rotation. Then use \eqref{eq:source-4} at that
successor, with its predecessor as $W$, to obtain a further error less
than $\gamma/2$. Combining coordinates gives
\eqref{eq:spectral-vector-approximation}. The coordinate count,
number of basis functions, their uniform bounds, and $B$ are all fixed
independently of $k$.

The earlier tower through $W$ is also a tower of factors of $Z$:
$Z$ contains $Y$ and hence every predecessor of $Y$. All measures,
actions, limit joins, and approximation properties of those factors
are unchanged. The universal earlier induction assertion therefore
applies in $Z$ to the map $b$. It says that the difference law of $b$
over $W$ is wild. Its Fourier coefficient at $a\in E$ is
\begin{equation}\label{eq:spectral-conditional-coefficient}
 p(a)=\left\|\mathbb E_Z(\psi_a(b)\mid W)\right\|_2^2\ge0,
\end{equation}
by conditional independence. Lemma~\ref{lem:line-average} applies to
this coefficient function.

\paragraph{Fixing the number of tests before choosing a small parameter.}
At any fixed fiber of $W$, the two coefficient arrays for two
approximations in \eqref{eq:spectral-vector-approximation} belong to
a fixed bounded subset of $\C^{2q}$. Partition this subset into a
finite number $L$ of bins so fine that, within one bin, the corresponding
approximating functions differ pointwise by at most $\tau$ for each
of the two vectors. For example, coordinate diameters less than
\[
 \frac{\tau}{1+\sum_{r=1}^q\|H_r\|_\infty}
\]
suffice. The number $L$ is independent of the fiber and of the chosen
rotations. Fix an integer $M\ge2$ so large that
\[
 \frac{0.9M}{L}>10m^2.
\]
All approximation data have been fixed before this choice.

There is a nonzero $t\in F$ with $|t|<1/M$ for which every value
\[
 p\bigl(n(j-l)t\bigr),\qquad
 1\le n\le N,\quad 1\le j,l\le M,\quad j\ne l,
\]
is as small as we prescribe. To see this, choose a nonnegative
$\phi\in C_c((-1/M,1/M))$ with positive integral. Each indicated
nonnegative test has zero $\phi$-weighted $I$-integral by
\eqref{eq:source-2}. Their finite sum does too. If that sum were
bounded below by a prescribed positive number at every point where
$\phi>0$, positivity of $I$ would contradict $I(\phi)>0$. Thus some
such point makes their sum, and hence every test, as small as required.
Choosing this bound below one also ensures $t\ne0$, since $p(0)=1$.

Set $x_j=jt$. Choose the prescribed smallness so that Markov's
inequality and a finite union bound give, on a set of $W$-probability
greater than $0.99$, all of the bounds
\[
 \left|\mathbb E_Z\bigl(\psi_{n(x_j-x_l)}(b)\mid W=w\bigr)\right|<c
 \qquad(1\le n\le N,\ j<l).
\]
More explicitly, the probability of failure of any one bound is at
most $p(n(x_j-x_l))/c^2$, by
\eqref{eq:spectral-conditional-coefficient}; there are finitely many
bounds. Apply \eqref{eq:spectral-small-arc} to the conditional law
of $\psi_{x_j-x_l}(b)$. On the same set of fibers,
\begin{equation}\label{eq:source-5}
 \mathbb P_Z\bigl(
   |\psi_{x_j}(b)-\psi_{x_l}(b)|\le4\rho\mid W=w
                  \bigr)<\frac1{4m^2}
 \qquad(1\le j<l\le M).
\end{equation}

\paragraph{A common coefficient bin on a good fiber.}
Each $|x_j|<1$ is a sum of two elements of $K$. One explicit choice is
\[
 k_j=\frac{x_j}{2}+i\sqrt{1-\frac{x_j^2}{4}},\qquad
 k'_j=\frac{x_j}{2}-i\sqrt{1-\frac{x_j^2}{4}}.
\]
The square roots belong to $F$, and $x_j=k_j+k'_j$. Apply the fixed
approximations \eqref{eq:spectral-vector-approximation} to
$U_{k_j}H$ and $U_{k'_j}H$ for every $j$.

Let $e_j^+(w),e_j^-(w)$ be the respective conditional squared
$L^2$ errors on the fiber of $Y$ over $w$. Their global integrals
are less than $\gamma^2$. If at least $0.1M$ indices fail to have
both conditional $L^2$ errors less than $\tau$, then
\[
 \frac1M\sum_{j=1}^M(e_j^+(w)+e_j^-(w))\ge0.1\tau^2.
\]
Markov's inequality and \eqref{eq:spectral-error-choice} bound the
probability of this event by
\[
 \frac{20\gamma^2}{\tau^2}<0.001.
\]
In particular, on a set of fibers of probability greater than $0.99$,
more than $0.9M$ indices have both conditional errors less than $\tau$.

Fix a fiber $w$ satisfying this property and all of
\eqref{eq:source-5}, and avoid the null sets for the finitely many
identities and bounds used here. Bin the two coefficient arrays at
this $w$ as above. There is a subset $J$ of the good indices in a
single bin with $|J|>10m^2$. Fix $j_0\in J$. For either type of
encoding the conditional $L^2$ distance from index $j$ to index $j_0$
is at most $3\tau$: the two approximation errors each cost less than
$\tau$, and the approximants differ pointwise by at most $\tau$.

For a uniform $j\in J$ and a conditional point $y$ over $w$, the
probability that either encoding distance is at least $s$ is at most
\[
 \frac{18\tau^2}{s^2}=0.0018.
\]
Thus, with probability greater than $0.95$, both phase lists for
index $j$ match their respective representative lists for $j_0$
within $\rho$. Let $q(y)$ be the fraction of indices in $J$ for
which both matches hold. We have
\begin{equation}\label{eq:spectral-matching-fraction}
 \mathbb E(q\mid W=w)>0.95.
\end{equation}

\paragraph{The collision contradiction.}
Fix such a fiber point $y$ and any point $b_0$ in the support of
$\beta_y$. For every matching index $j$, the two phases
$\psi_{k_j}(b_0)$ and $\psi_{k'_j}(b_0)$ are within $\rho$ of entries
of their respective representative lists. Consequently
\[
 \psi_{x_j}(b_0)=\psi_{k_j}(b_0)\psi_{k'_j}(b_0)
\]
is within $2\rho$ of one of the at most $m^2$ products of entries
from those two lists. The matchings need not identify the same center
in the two lists; allowing all $m^2$ products accounts for this.

Assign each of the $q(y)|J|$ matching indices to one such product.
By the sum-of-squares inequality, the number of ordered pairs assigned
to the same product is at least $q(y)^2|J|^2/m^2$. Removing diagonal
pairs discards at most $|J|$ pairs. Hence, for independent uniform
$j,l\in J$,
\[
 \mathbb P_{j,l}\bigl(j\ne l,\ 
 |\psi_{x_j}(b_0)-\psi_{x_l}(b_0)|\le4\rho\bigr)
 \ge\frac{q(y)^2}{m^2}-\frac1{|J|}.
\]
This bound holds for every possible center $b_0$, so it can be
integrated first over $\beta_y$ and then over the conditional law of
$y$ given $w$. Jensen's inequality and
\eqref{eq:spectral-matching-fraction} give a lower bound strictly
greater than
\[
 \frac{0.95^2}{m^2}-\frac1{|J|}
 >\frac{0.95^2-0.1}{m^2}
 >\frac1{4m^2}.
\]
But \eqref{eq:source-5} bounds the conditional collision probability
for every distinct pair in the original set of $M$ indices by
$1/(4m^2)$. Averaging those bounds over the selected $J$ gives an
upper bound smaller than $1/(4m^2)$. Choosing $J$ after fixing $w$
is harmless because all those pair bounds hold simultaneously at that
fiber. This contradiction completes the induction step, and hence the
universal transfinite induction.
\end{proof}
